% !PRINTED 469
\weberpage{469}

and there are thus eight numerical coefficients to compute, which
can be found by comparing a few terms in the expressions of the
$a_\nu$ in terms of the roots on the one hand, and in terms of the invariants (\S.\,123)
on the other, and which, besides rational numbers, can contain only
$\sqrt{-7}$.

We shall at first compute these coefficients only under the
assumption that $f=0$ ${}^1$).

One then needs only to compute the first terms (with $x_1^5x_2$, $x_1^{10}x_2^2$)
in the power sums $\Sigma z_r^3$ and $\Sigma z_r^6$ and to apply the
\gesperrt{Newton} formulae. The last coefficient results
from the one term $x_1^{14}$ in the product of the $z_r$. One
finds first
\[
\Sigma z_r^3 = -3\cdot7\cdot\frac{1-\sqrt{-7}}{2}\,x_1^5x_2\dots
\]
\[
\Sigma z_r^6 = \left[6\cdot7 + 15\cdot7\left(\frac{1-\sqrt{-7}}{2}\right)^2\right]x_1^{10}x_2^2\dots
\]
and from these the sought resolvent
\weq{9}{z^7 - 7\cdot\frac{1-\sqrt{-7}}{2}\,\Delta z^4 - 7\cdot\frac{5+\sqrt{-7}}{2}\,\Delta^2 z - C = 0.}
If one wishes to reduce this equation to another which
depends only on the ratios of the $x_1, x_2, x_3$, one sets
\weq{10}{z = \frac{Cu}{\Delta^2}, \quad g = \frac{\Delta^7}{C^3},}
whereby one obtains from (9) an equation in which only the
one parameter $g$ still occurs, namely
\weq{11}{u^7 - 7\,\frac{1-\sqrt{-7}}{2}\,g\,u^4 - 7\,\frac{5+\sqrt{-7}}{2}\,g^2\,u - g^2 = 0.}
If one makes the same substitution in the general resolvent,
in which $f$ is not set $=0$, one has to introduce a
further parameter
\weq{12}{h = \frac{f\Delta^4}{C^2}}
as well, and the coefficients of the resolvent become, apart from
%CONT>>

\par\noindent\rule{0.3\linewidth}{0.4pt}\par
${}^1$) This case is of particular interest because, like
the icosahedral equation, it leads to the transformation equations from the
theory of elliptic functions.
